\documentclass[Main.tex]{subfiles}
\newcommand{\underwrite}[3][]{% \underwrite[<thickness>]{<numerator>}{<denominator>}
  \genfrac{}{}{#1}{}{\textstyle #2}{\scriptstyle #3}
}

\renewcommand{\thesection}{\arabic{section}}

\begin{document}
\begin{titlepage}
\thispagestyle{empty}
\begin{center}

{\textbf{\large{HANOI UNIVERSITY OF SCIENCE AND TECHNOLOGY }}}\\[4cm]

{\textbf{\huge{GRADUATION THESIS}}}\\[1cm]
% Working title; to be revisited once the November runs are in.
{\textbf{\Large{Optimizing Universal Photometric Stereo via Wavelet-Energy Saliency and Latent Carrier Tokens}}\\[1cm]

{\textbf{\large{LÊ NHẬT QUANG}}}\\
{\large{quang.ln225522@sis.hust.edu.vn}}\\[0.5cm]

{\textbf{\large{Program: Data Science and Artificial Intelligence}}}\\

\vspace{4cm}
\begin{table}[H]
\centering
\resizebox{\textwidth}{!}{%
\begin{tabular}{ll}
 \multicolumn{1}{l}{\textbf{Supervisor:}} & Ngô Thành Trung \hspace{0.5cm}
  \\[0.5cm]
 \multicolumn{1}{l}{\textbf{Department:}} & Computer Engineering \hspace{0.5cm} \\ [0.5cm]
 \multicolumn{1}{l}{\textbf{School:}} & School of Information and Communications Technology \hspace{0.5cm} \\ [5.5cm]
\multicolumn{2}{c}{\textbf{HANOI, 11/2026}}
\end{tabular}%
}
\end{table}}
\end{center}

\end{titlepage}

% =====================================================================
\chapter*{Thesis Outline Proposal}

This document proposes the structure of the thesis for the supervisor's review. The thesis is an
iterative investigation of modifications to SDM-UniPS in which each result motivated the next
step, so its core is organized as one chapter per modification rather than as a single
methodology chapter followed by a single results chapter. The template's research-thesis
structure is kept where it applies: the Introduction uses the template's sections, the Literature
Review respects its length limit, the chapter on the baseline and experimental framework plays
the role of the template's Methodology and Simulation Method, and every modification chapter
follows the same internal order (method, then results, then analysis), with an overview at the
start of each chapter and a conclusion at its end.

\section*{Research questions}

\begin{enumerate}[label=\textbf{RQ\arabic*.}, leftmargin=*]
  \item Does a wavelet-convolution backbone improve SDM-UniPS?
  \item Does wavelet-energy saliency sampling improve SDM-UniPS, and if not, why?
  \item Does a carrier-token / deformable-attention decoder improve SDM-UniPS, and does it remove
    the penalty of non-uniform sampling?
  \item Once the sample-set bias is removed by decoupling the context set from the query set, can
    non-uniform sampling improve accuracy?
\end{enumerate}

% =====================================================================
\section*{Front matter}
Title page, declaration of originality, acknowledgement (100--150 words), abstract (200--350
words, one paragraph, written last), table of contents, list of figures, list of tables, list of
abbreviations.

% =====================================================================
\section*{Chapter 1. Introduction \normalfont{(3--6 pages)}}
\begin{itemize}
  \item[1.1] \textbf{Problem Statement.} Surface-normal recovery under unknown, spatially varying
    lighting (universal photometric stereo), and why it matters.
  \item[1.2] \textbf{Background and Problems of Research.} Current methods up to SDM-UniPS, and the
    three limitations targeted: (i) large depthwise kernels low-pass the shading gradients normals
    are read from, (ii) training pixels are drawn uniformly regardless of where information lies,
    and (iii) the decoder's context comes only from the sampled pixels.
  \item[1.3] \textbf{Research Objectives and Conceptual Framework.} Scope (normals only, fixed
    camera, trained on synthetic data, evaluated on DiLiGenT), RQ1--RQ4, and the modification
    proposed for each limitation.
  \item[1.4] \textbf{Contributions.}
  \item[1.5] \textbf{Organization of Thesis.}
\end{itemize}

% =====================================================================
\section*{Chapter 2. Background and Literature Review \normalfont{(at most 10 pages)}}
\begin{itemize}
  \item[2.1] \textbf{Background.} Only concepts later chapters use: image formation (radiometry,
    BRDF, Lambertian model); photometric stereo settings (calibrated, uncalibrated, universal),
    normal maps and angular error; building blocks (ConvNeXt, UPerNet, attention, Haar discrete
    wavelet transform, deformable attention, importance sampling).
  \item[2.2] \textbf{Related works}, by topic: classical photometric stereo (condensed from
    Project III); learning-based calibrated photometric stereo; uncalibrated and universal
    photometric stereo, ending at SDM-UniPS; frequency and wavelet methods in vision (WTConv,
    FUPS-Net); non-uniform and importance sampling in training; efficient attention (deformable
    attention, carrier tokens).
  \item[2.3] \textbf{Summary and research gap.}
\end{itemize}

% =====================================================================
\section*{Chapter 3. Baseline and Experimental Framework}
\begin{itemize}
  \item[3.1] \textbf{SDM-UniPS walkthrough}, with one figure marking every modification point.
  \item[3.2] \textbf{Training pipeline} for the inference-only release: loss and the uniform
    2,048-pixel training draw, augmentation, per-image normalization, optimization (AdamW,
    warmup--stable--decay schedule), and deviations from the released code.
  \item[3.3] \textbf{Baseline correction.} The dead coarse branch of the fusion head and the
    GroupNorm fix, which together define Model A.
  \item[3.4] \textbf{Datasets.} hdlong-complexv1, PolarPS (synthesis of the training renders,
    scene-level splits) and DiLiGenT.
  \item[3.5] \textbf{Evaluation protocol.} Fixed evaluation pixels and renders, so every variant is
    evaluated on identical data; DiLiGenT at $K = 4, 8, 16$ images with 10 paired trials per
    object; the noise reference from two trainings of the same model.
  \item[3.6] \textbf{Implementation details.} Hardware, schedule, hyperparameters.
\end{itemize}

% =====================================================================
\section*{Chapter 4. B1: Wavelet Convolution Backbone \normalfont{(RQ1)}}
Motivation; WTConv as a drop-in replacement for the $7\times7$ depthwise convolution; results
against Model A, reproduced by a second training of B1.

% =====================================================================
\section*{Chapter 5. B2: Wavelet-Energy Saliency Sampling \normalfont{(RQ2)}}
\begin{itemize}
  \item[5.1] \textbf{Premise validation.} Does the energy map point at surface curvature
    (Phase 0, with a shuffled-energy control)?
  \item[5.2] \textbf{Sampler design.} Energy map, exclusion of the silhouette band, mixture with
    the uniform draw.
  \item[5.3] \textbf{Analysis of the sampling objective.} (i) Every masked pixel keeps a
    probability of at least $\lambda / n$; (ii) the mismatch between training and validation: the
    unweighted sampler minimizes $\sum_i p_i \ell_i$ while validation measures
    $\frac{1}{n}\sum_i \ell_i$; (iii) importance weights $w_i = 1/(n p_i)$: expectation 1 under the
    draw, maximum $1/\lambda$, and unbiased when a pixel's loss does not depend on the other pixels
    drawn with it; (iv) the sample-set bias when that condition fails. Items (ii) and (iv) are the
    two hypotheses the diagnosis tests. Full derivations in an appendix.
  \item[5.4] \textbf{Temperature and mixture sweep.}
  \item[5.5] \textbf{Diagnosis.} In order: the focus-on-creases claim is not supported
    (curvature deciles); the shifted training target is ruled out (importance-weighted run);
    variance reduction is refuted; a sample-set bias is measured.
\end{itemize}

% =====================================================================
\section*{Chapter 6. C: Carrier Tokens and Deformable Attention \normalfont{(RQ3)}}
Motivation from the sample-set bias of Chapter 5; design and cost; C with uniform sampling against
B1; the $2\times2$ decoder $\times$ sampler comparison with its pre-registered verdict.

% =====================================================================
\section*{Chapter 7. Decoupled Context/Query Decoder \normalfont{(RQ4, November 2026)}}
\begin{itemize}
  \item[7.1] \textbf{Design.} Built on B1: a uniform context set attends within itself as in B1;
    each query attends to the context set only. Why this makes the importance-weighted loss
    unbiased again (builds on Section 5.3).
  \item[7.2] \textbf{Run Z} (WESS with importance weights). Against B2-IW, which differs only in
    the decoder; against B1.
  \item[7.3] \textbf{Run Y} (learned loss-prediction sampler). Against Z, which differs only in
    the sampler. Conditional on time; otherwise moved to future work.
\end{itemize}

% =====================================================================
\section*{Chapter 8. Discussion}
All variants in one table; cost against gain; the gap between synthetic and real data (why models
that win on synthetic held-out data do not win on DiLiGenT at high $K$); threats to validity.

% =====================================================================
\section*{Chapter 9. Conclusions}
\begin{itemize}
  \item[9.1] \textbf{Summary.} Answers to RQ1--RQ4 and limitations.
  \item[9.2] \textbf{Suggestion for Future Works.}
\end{itemize}

\section*{References and Appendices}
References in IEEE style. Appendices: full hyperparameter tables, per-object DiLiGenT results,
derivations for Section 5.3, probe protocols, training infrastructure (resuming, warm-starting, checkpoint backup).

% =====================================================================
\section*{Status of experiments}
\begin{table}[H]
\centering
\begin{tabular}{lll}
\hline
\textbf{Run} & \textbf{Content} & \textbf{Status} \\
\hline
A        & SDM-UniPS with GroupNorm fusion head            & done \\
B1       & WTConv backbone (and a second training)         & done \\
B2       & WESS, temperature and mixture sweep             & done \\
B2-IW    & WESS with importance weights                    & done \\
C        & Carrier tokens + deformable attention, uniform  & done \\
C + WESS & Carrier tokens + deformable attention, WESS     & done \\
Z        & Decoupled decoder (on B1), WESS with importance weights & November 2026 \\
Y        & Decoupled decoder (on B1), learned sampler      & November 2026, conditional \\
\hline
\end{tabular}
\end{table}

\section*{Changes from the initial roadmap}
\begin{itemize}
  \item The WESS draw is a single mixture of a uniform and an energy-based distribution, with the
    silhouette band kept at the uniform rate, instead of a 512 + 1,536 split of the budget.
  \item Model C has no repulsion loss.
  \item The decoupled decoder and the learned sampler were not in the roadmap; they follow from the
    sample-set bias measured in the WESS diagnosis.
\end{itemize}

\end{document}
